\paragraph{Close the recorded balances under explicit carry policies.}
The period report leaves profit in the temporary income and expense accounts.
\texttt{CompanyBookCarry/1} transfers those recorded balances to a named equity
account and constructs the next provisional opening. The caller identifies that
account explicitly; an account label cannot select it. Each temporary account
receives the negative of its closing balance, while the selected equity account
receives their sum. With debit-positive storage, profit therefore produces an
additional equity credit and a loss an equity debit. The journal sums to zero,
and cash and the other permanent accounts remain unchanged.

In the preceding company example, profit of USD~25 is transferred to equity,
leaving USD~130 of opening equity for the next period and zero income and expense
balances. An otherwise analogous USD~3 loss reduces equity by USD~3. If the
company names a separate retained-income account, the contribution account remains
unchanged. Zero profit produces no net retained-equity transfer, although offsetting
income and expense balances still require closing entries. The next opening report
has no current-period profit and reconciles to the preceding closing equity.

The method recomputes the period report and requires the closing snapshot to be
dated the last included day before the period's exclusive endpoint. The next
period starts at that endpoint and ends at an explicitly supplied later date;
calendar years are not assumed. This date rule does not establish that the event
record is complete. The method carries recorded balances provisionally and keeps
fiscal-completeness and accounting-admission flags false.

Every typed memory leaf needs its own declared policy: preserve its value, reset
it to an explicitly typed literal, or, for a date leaf, set it to the next period's
start. Each policy records its rationale. Missing policies and type mismatches
prevent carry. The engine does not silently reset payment counters, advance tax
clocks or infer an opening tax basis. Such choices can be type-correct yet
substantively wrong, so their provenance remains part of the proposed result.
The existing scalar memory interpreter applies the declared replacements.

The next book clears prior-period cash flows but retains lifetime event identities
and adds a new close identity. A returned archive record links the prior book,
report, closing journal result, next opening and memory policies by digest.
Application recomputes the full proposal. This is a returned provisional history
record, not a durable storage commit. Comparisons reproduce the retained native
carry's balances, period and imported lot memory; its other tax, interest and
archive conventions remain separate native rules. Consecutive-period tests retain
earlier event and close identities rather than making them available for reuse.
